
Notons aussi que, \(X\) étant régulier donc normal, ses composantes connexes sont identiques à ses composantes irréductibles, lesquelles dominent \(S\), de sorte que l'homomorphisme \(H^0(X,\mathbb Z)\to H^0(X_t,Z)\) induit par l'inclusion \(X_t\to X\) est bijectif. Par suite l'homomorphisme \(H^0(X,\mathbb Z)\to H^0(X_s,\mathbb Z)\) induit par l'inclusion \(X_s\to X\) peut aussi s'interpréter comme un homomorphisme
\begin{equation}\tag{7.7.3}
  H^0(X_t,\mathbb Z)\longrightarrow H^0(X_s,\mathbb Z),
\end{equation}
qui mérite encore le nom de «~homomorphisme de spécialisation~». Ceci posé, on vérifie immédiatement sur les définitions que l'homomorphisme (7.7.1) est compatible avec les structures d'\emph{algèbres augmentées} sur \(H^0(X_t,\mathbb Z)\) resp. \(H^0(X_s,\mathbb Z)\) des deux membres définies dans (IV 2.8), compte tenu de l'homomorphisme précédent (7.7.3).

Supposant toujours \(X\) séparé, ce qui précède implique en particulier que (7.7.1) induit un homomorphisme d'anneaux gradués
\begin{equation}\tag{7.7.4}
  \sigma^\bullet:\Gr^\bullet(X_t)\longrightarrow\Gr^\bullet(X_s),
\end{equation}
et utilisant la notion de classe de Chern de (V 6), on trouve un diagramme commutatif
\[
\begin{tikzcd}[column sep=large,row sep=large]
K^\bullet(X_t) \arrow[r,"c_{X_t}"] \arrow[d] &
\Ch\,\Gr^\bullet(X_t) \arrow[d]\\
K^\bullet(X_s) \arrow[r,"c_{X_s}"] &
\Ch\,\Gr^\bullet(X_s).
\end{tikzcd}
\]

\subsection*{7.8. Spécialisation pour \(\Pic\).}
Gardons les hypothèses de 7.7 : \(X\) séparé, \(X_t\) et \(X_s\) réguliers. Tenant compte de l'isomorphisme canonique (5.3.2) :
\[
  \Gr^1\simeq \Pic ,
\]
l'homomorphisme (7.7.4) nous donne donc un homomorphisme canonique
\begin{equation}\tag{7.8.1}
  \Pic(X_t)\longrightarrow \Pic(X_s).
\end{equation}
